\section{Dataset and Software Framework}

\subsection{ROOT input files}

The project uses the CERN Open Data record \emph{Matter Antimatter Differences
(B meson decays to three hadrons)}~\cite{b2hhhdata}, released by the LHCb
collaboration in 2017 under a CC0 licence. The record contains three ROOT files
with reconstructed $B$-candidate ntuples. The two data files,
\begin{center}
  \texttt{data/raw/B2HHH\_MagnetDown.root}, \qquad
  \texttt{data/raw/B2HHH\_MagnetUp.root},
\end{center}
hold the same data-taking period recorded with the two polarities of the LHCb
dipole magnet; the third, \texttt{PhaseSpaceSimulation.root}, is kept for the
extensions discussed in Sec.~\ref{sec:conclusion}. Both polarities are used for
the raw asymmetry. Reversing the magnetic field swaps the side of the detector
towards which each charge is bent, so averaging the two polarities cancels
left--right detection asymmetries to first order
(Sec.~\ref{sec:cp}); MagnetUp also provides an independent sample on which a
feature seen in MagnetDown can be checked. The selection, the classifier study and
the exploratory notebooks were developed on MagnetDown.

The data are stored in the ROOT format~\cite{root}, a standard file and object format in high-energy physics. The relevant event information is contained in a tree named \texttt{DecayTree}. A ROOT \texttt{TTree} stores columnar data: each branch corresponds to one observable, such as a momentum component, charge, PID response, vertex quality, or flight distance. This format is efficient for large event samples because only the branches needed for a given analysis step need to be read.

\subsection{Python analysis workflow}

The analysis is implemented as a Python workflow. Reusable functions live under \texttt{src/}; exploratory notebooks under \texttt{notebooks/}; the fit entry points under \texttt{scripts/}; and the toy-MC, systematic, charge-shape and Dalitz studies under \texttt{studies/}. Every study writes its numerical output as JSON and CSV under \texttt{results/}, so each number quoted in this report can be traced to the file that produced it. The main software components are listed in Table~\ref{tab:software}; they are drawn from the scientific Python stack~\cite{numpy,matplotlib} and from the Scikit-HEP ecosystem~\cite{scikithep}.

\begin{table}[H]
  \centering
  \caption{Main software tools used in the analysis workflow.}
  \label{tab:software}
  \begin{tabular}{ll}
    \toprule
    Tool & Role in the analysis \\
    \midrule
    ROOT & Standard HEP file format and ntuple storage \\
    \texttt{uproot} & Reading ROOT TTrees directly in Python \\
    \texttt{pandas} & Tabular event manipulation and selections \\
    \texttt{numpy} & Vectorised numerical calculations \\
    \texttt{scipy} & Special functions used in the mass models \\
    \texttt{iminuit} & \textsc{Minuit} interface for the likelihood minimisation \\
    \texttt{matplotlib} & Publication-style figures \\
    \texttt{scikit-learn} & Random Forest classifier and performance metrics \\
    \bottomrule
  \end{tabular}
\end{table}

The branch set used throughout the analysis includes the momentum components \texttt{H1\_PX}, \texttt{H1\_PY}, \texttt{H1\_PZ} through \texttt{H3\_PZ}, the three track charges, PID variables such as \texttt{H1\_ProbK}, and B-candidate quality variables such as \texttt{B\_VertexChi2} and \texttt{B\_FlightDistance}. A shared selection helper loads these branches, reconstructs the B-candidate mass, and applies the common baseline cuts. This avoids duplicated logic between the mass, machine-learning, CP-asymmetry, and Dalitz notebooks.

The workflow is intentionally modular. The main reusable components are the invariant-mass reconstruction module \path{src/reconstruction/invariant_mass.py}, the mass models and likelihood fit in \path{src/fitting/}, the shared selection in \path{src/physics/selection.py}, the Dalitz-kinematics module \path{src/physics/dalitz.py}, the plotting helpers in \path{src/plotting/mass_plots.py}, and the classifier utilities in \path{src/ml/signal_background.py}. This separation mirrors the structure of a maintainable HEP analysis, where physics definitions, fit models, plotting conventions and machine-learning utilities are reusable across studies rather than duplicated in each of them.
